% ============================================================
%  PAPER 2: Medical Image Analysis Format
%  Title: Systematic evaluation of twelve bandpass filters for
%         FFT-based segmentation of phase-contrast microscopy
%         across eight cell lines and thirteen degradation types
% ============================================================
%  Target: Medical Image Analysis (Elsevier)
%  Author: Prof. Dr. Md. Enamul Hoque
%  Institution: Shahjalal University of Science and Technology
% ============================================================

\documentclass[11pt,a4paper]{article}

% ---- Packages ----
\usepackage[margin=2.5cm]{geometry}
\usepackage{graphicx}
\usepackage{booktabs}
\usepackage{amsmath,amssymb}
\usepackage{siunitx}
\usepackage{caption}
\usepackage{subcaption}
\usepackage{hyperref}
\usepackage{xcolor}
\usepackage{float}
\usepackage{setspace}
\usepackage{parskip}
\usepackage{multirow}
\usepackage{colortbl}
\usepackage{tikz}
\usepackage{pgfplots}
\usepackage{enumitem}
\usepackage{wrapfig}
\usepackage{makecell}
\usepackage{xspace}

\usetikzlibrary{shapes,arrows,positioning,calc,decorations.pathreplacing,patterns,fit,backgrounds}
\pgfplotsset{compat=1.18}

\hypersetup{
    colorlinks=true,
    linkcolor=blue!70!black,
    citecolor=blue!70!black,
    urlcolor=blue!70!black
}

\onehalfspacing

% ---- Custom Commands ----
\newcommand{\eg}{e.g.,\xspace}
\newcommand{\ie}{i.e.,\xspace}
\newcommand{\etal}{et~al.\xspace}
\newcommand{\fft}{\textsc{fft}\xspace}
\newcommand{\iou}{\textsc{iou}\xspace}
\newcommand{\svm}{\textsc{svm}\xspace}
\newcommand{\rf}{\textsc{rf}\xspace}
\newcommand{\lr}{\textsc{lr}\xspace}
\newcommand{\pca}{\textsc{pca}\xspace}
\newcommand{\bce}{\textsc{bce}\xspace}
\newcommand{\nii}{\textsc{n2v}\xspace}
\newcommand{\care}{\textsc{care}\xspace}
\newcommand{\debcr}{\textsc{debcr}\xspace}
\newcommand{\piddpm}{\textsc{pi-dddpm}\xspace}
\newcommand{\dog}{\textsc{dog}\xspace}
\newcommand{\lo}{\textsc{lq}\xspace}
\newcommand{\hq}{\textsc{hq}\xspace}

% ---- Title ----
\title{\textbf{Systematic evaluation of twelve bandpass filters for FFT-based segmentation of phase-contrast microscopy across eight cell lines and thirteen degradation types}}

\author{
Prof.~Dr.~Md.~Enamul Hoque\thanks{Corresponding author: \texttt{enamul.physics@sust.edu}} \\
\textit{Department of Physics, Shahjalal University of Science and Technology} \\
Sylhet, Bangladesh
}

\date{\today}

% ============================================================
\begin{document}
\maketitle

% ============================================================
%  ABSTRACT (150-250 words, structured)
% ============================================================
\begin{abstract}
\textbf{Background and Objective:} Phase-contrast microscopy is widely used for label-free cell imaging, and frequency-domain bandpass filtering is the standard preprocessing step for segmentation. However, no systematic evaluation of filter types across cell lines and image quality levels has been performed. Different studies use different filters with no justification, and the ``best'' filter remains unknown.

\textbf{Methods:} We evaluated 12 bandpass filter types (Ideal, Butterworth, Gaussian, Chebyshev I/II, Elliptic, Laplacian-BP, Homomorphic, Gabor, Difference-of-Gaussians, Trapezoidal, Cosine-tapered, and Parametric) across 8 cell lines from the LIVECell dataset, 4 degradation types (Gaussian noise, defocus blur, illumination shading, combined), and 25 blur levels from the BBBC005 dataset. In total, over 20,000 segmentation evaluations were performed using Otsu thresholding and U-Net segmentation with 5-fold cross-validation.

\textbf{Results:} No single filter is universally optimal. DoG achieves the best segmentation for 3/8 cell lines, Homomorphic for 2/8, and Butterworth/Elliptic/Gaussian each for 1/8. Filter improvements on low-quality images are 10--100$\times$ smaller than on high-quality images. Filter transfer efficiency from high-quality to low-quality images is consistently below 15\%. Adaptive, cell-line-specific filter selection improves mean intersection-over-union by $+0.130$ over fixed filtering.

\textbf{Conclusions:} We provide the first comprehensive, evidence-based filter selection guide for phase-contrast microscopy. The adaptive selection algorithm and quality-aware recommendations are publicly available.

\textbf{Keywords}: Bandpass filtering, phase-contrast microscopy, cell segmentation, Fourier transform, LIVECell, image quality, adaptive filtering
\end{abstract}

% ============================================================
%  1. INTRODUCTION (~600 words)
% ============================================================
\section{Introduction}
\label{sec:intro}

Phase-contrast microscopy is one of the most widely used techniques for imaging transparent biological specimens without staining~\cite{zernike1942}. The resulting images encode structural information through intensity variations caused by refractive index differences between cellular components and the surrounding medium. Automated segmentation of cells in phase-contrast images is essential for high-throughput screening, drug discovery, and cell biology research~\cite{edlund2021livecell}.

Frequency-domain bandpass filtering is the standard preprocessing step for phase-contrast image segmentation~\cite{gonzalez2018digital}. By selectively preserving spatial frequencies corresponding to cell boundaries while suppressing low-frequency illumination gradients and high-frequency noise, bandpass filters can significantly improve segmentation accuracy. However, the literature lacks consensus on which filter type to use. Different studies apply Ideal~\cite{castleman1979}, Butterworth~\cite{lloyd1993}, Gaussian~\cite{bray2012}, or Homomorphic filters with no systematic justification. This is problematic because: (1) different filter types have fundamentally different frequency response characteristics (sharp vs.\ smooth transition, passband ripple, stopband attenuation); (2) filter performance may depend on cell line morphology; and (3) the interaction between image quality and filter performance is poorly understood.

Previous work on filter evaluation has been limited in scope. Most studies compare 2--3 filter types on a single cell line or a small dataset~\cite{castleman1979, gonzalez2018digital}. No study has systematically evaluated filter performance across multiple cell lines with ground-truth annotations, multiple degradation types, and controlled blur levels. Furthermore, the question of whether filter parameters optimized on high-quality (HQ) images transfer to low-quality (LQ) images---a critical question for automated screening---has not been addressed.

Here we present the first comprehensive evaluation of 12 bandpass filter types for phase-contrast microscopy segmentation. Using the LIVECell dataset (3,727 images, 8 cell lines, 258,569 annotated instances) and the BBBC005 dataset (768 images, 25 blur levels), we perform over 20,000 segmentation evaluations across 4 degradation types and 25 blur levels. Our contributions are: (1) a systematic comparison of 12 filter types across 8 cell lines; (2) quantification of the quality-dependent performance gap; (3) measurement of filter transfer efficiency from HQ to LQ images; (4) an adaptive filter selection algorithm; and (5) evidence-based, degradation-specific filter recommendations.

% ============================================================
%  2. RELATED WORK (~400 words)
% ============================================================
\section{Related Work}
\label{sec:related}

\textbf{Fourier analysis in cell biology.} The fast Fourier transform has been applied to cell biology for texture classification~\cite{castleman1979}, cell cycle analysis~\cite{lloyd1993}, and morphological profiling~\cite{bray2012}. The power spectral density encodes cell size (low-frequency), cell density (total power), and morphological regularity (spectral shape). However, systematic evaluation of FFT features across multiple cell lines with ground-truth annotations has been limited.

\textbf{Bandpass filtering for segmentation.} Frequency-domain filtering is a standard preprocessing step in image processing~\cite{gonzalez2018digital}. Butterworth filters are recommended for general-purpose applications due to their maximally flat passband, while Gaussian filters are preferred when spatial-domain artifacts (ringing) must be avoided. Homomorphic filtering is specifically designed for multiplicative noise and illumination correction. Despite these well-known properties, no study has systematically compared filter types for cell segmentation.

\textbf{Filter selection in microscopy.} In fluorescence microscopy, filter selection is guided by the known emission spectrum. In phase-contrast microscopy, the situation is more complex because the contrast transfer function produces characteristic halo artifacts that dominate specific frequency bands. The optimal filter must preserve cell boundary frequencies while suppressing halo artifacts, and the optimal cutoff frequencies likely depend on cell line morphology.

\textbf{Image quality and filter performance.} Image quality degradation in microscopy arises from multiple sources: photon noise, defocus, illumination non-uniformity, and stage drift. The effect of these degradations on frequency-domain filter performance has not been systematically studied.

% ============================================================
%  3. METHODS (~800 words)
% ============================================================
\section{Methods}
\label{sec:methods}

\subsection{Dataset}
The LIVECell 2021 dataset~\cite{edlund2021livecell} contains 3,727 phase-contrast images (704$\times$520 px) from 8 cell lines with 258,569 annotated cell instances (808 annotated images). The BBBC005 dataset~\cite{ljosa2012annotated} contains 768 images at 25 controlled blur levels ($s_{01}$--$s_{25}$).

\subsection{Synthetic Degradation Pipeline}
Four degradation types were applied: (D1) Gaussian noise ($\sigma{=}50$); (D2) Gaussian blur ($\sigma{=}4$); (D3) multiplicative illumination shading ($\alpha{=}0.5$); (D4) combined mild. Total: 16,912 synthetic LQ images.

\subsection{Bandpass Filter Library}
12 filter types were implemented in the frequency domain:
\begin{equation}
I_{\text{filt}}(x,y) = \mathcal{F}^{-1}\left[F(u,v) \cdot H_{\text{BP}}(u,v)\right] + \bar{I}
\end{equation}
The filters span: sharp-cutoff (Ideal, Elliptic), smooth-transition (Butterworth, Gaussian, Cosine-tapered), equiripple (Chebyshev I/II), specialized (Homomorphic, DoG, Gabor, Laplacian-BP), and hybrid (Trapezoidal, Parametric). Complete mathematical formulations are in Supplementary Table~S1.

\subsection{Segmentation}
Two segmentation methods were used: (1) Otsu thresholding (classical baseline); (2) U-Net~\cite{ronneberger2015u} with 5-fold stratified cross-validation. U-Net architecture: 4 encoding levels (32--64--128--256 features), combined BCE+Dice loss.

\subsection{Evaluation Metrics}
Primary metric: Intersection-over-Union (IoU). Secondary: Dice coefficient, precision, recall. Statistical significance: paired $t$-test with Bonferroni correction, Cohen's $d$ effect sizes.

\subsection{Adaptive Filter Selection}
Grid search over filter parameters for each cell line and degradation type. Parameter spaces: Butterworth $n \in \{1,2,3,4\}$, $d_{\text{low}} \in [0.005, 0.05]$, $d_{\text{high}} \in [0.15, 0.50]$; DoG $\sigma_1 \in [0.01, 0.08]$, $\sigma_2 \in [0.05, 0.40]$; Homomorphic $\gamma_L \in [0.2, 0.7]$, $\gamma_H \in [1.5, 3.0]$. Best configuration selected by mean IoU on HQ annotated subset.

\subsection{Code Availability}
All code: \url{https://github.com/mjonyh/microscopic_images}.

% ============================================================
%  4. RESULTS (~1,200 words)
% ============================================================
\section{Results}
\label{sec:results}

% ---- Figure 1: Filter Taxonomy ----
\begin{figure}[H]
\centering
% TikZ schematic: Filter Taxonomy (standalone compile)
\documentclass[tikz,border=10pt]{standalone}
\usepackage{tikz}
\usetikzlibrary{positioning,arrows.meta}

\begin{document}
\begin{tikzpicture}[
    node distance=0.8cm and 1.5cm,
    cat/.style={rectangle, draw, fill=blue!5, minimum width=3.5cm, minimum height=0.7cm, align=center, font=\small},
    filt/.style={rectangle, draw, fill=white, minimum width=2.8cm, minimum height=0.6cm, align=center, font=\footnotesize},
    rec/.style={rectangle, draw, fill=green!15, minimum width=2.8cm, minimum height=0.6cm, align=center, font=\footnotesize},
    warn/.style={rectangle, draw, fill=yellow!15, minimum width=2.8cm, minimum height=0.6cm, align=center, font=\footnotesize},
    bad/.style={rectangle, draw, fill=red!10, minimum width=2.8cm, minimum height=0.6cm, align=center, font=\footnotesize},
    arrow/.style={-{Stealth[length=2mm]}, thick}
]

% Categories
\node[cat] (sharp) at (0,4) {\textbf{Sharp Cutoff}};
\node[cat] (smooth) at (5,4) {\textbf{Smooth Transition}};
\node[cat] (special) at (10,4) {\textbf{Specialized}};

% Sharp cutoff filters
\node[filt] (ideal) at (0,2.5) {Ideal\\(brick-wall)};
\node[warn] (cheb1) at (-1.5,1.2) {Chebyshev I\\(equiripple passband)};
\node[warn] (cheb2) at (1.5,1.2) {Chebyshev II\\(equiripple stopband)};
\node[bad] (elliptic) at (0,-0.1) {Elliptic\\(Cauer)};

% Smooth transition filters
\node[rec] (butter) at (5,2.5) {Butterworth\\(maximally flat)};
\node[rec] (gauss) at (5,1.2) {Gaussian\\(zero ringing)};
\node[rec] (cosine) at (5,-0.1) {Cosine-tapered\\(Hann)};
\node[filt] (trap) at (5,-1.4) {Trapezoidal\\(linear ramp)};

% Specialized filters
\node[rec] (dog) at (10,2.5) {DoG\\(blob detection)};
\node[rec] (homo) at (10,1.2) {Homomorphic\\(illumin. correction)};
\node[filt] (gabor) at (10,-0.1) {Gabor\\(orientation-selective)};
\node[filt] (lapl) at (10,-1.4) {Laplacian-BP\\(edge enhancement)};

% Arrows from categories
\draw[arrow] (sharp) -- (ideal);
\draw[arrow] (smooth) -- (butter);
\draw[arrow] (special) -- (dog);

% Legend
\node[draw, fill=green!15, minimum width=1.5cm, minimum height=0.5cm, font=\footnotesize] (leg1) at (1.5,-3) {Recommended};
\node[draw, fill=yellow!15, minimum width=1.5cm, minimum height=0.5cm, font=\footnotesize, right=0.3cm of leg1] (leg2) {Conditional};
\node[draw, fill=red!10, minimum width=1.5cm, minimum height=0.5cm, font=\footnotesize, right=0.3cm of leg2] (leg3) {Not recommended};

\end{tikzpicture}
\end{document}

\caption{\textbf{Bandpass filter taxonomy.} 12 filter types organized by transition sharpness and artifact characteristics. Color coding: green = recommended, yellow = conditional, red = not recommended for phase-contrast microscopy.}
\label{fig:taxonomy}
\end{figure}

% ---- 4.1 Filter Performance on HQ Images ----
\subsection{No single filter is universally optimal on high-quality images}

Table~\ref{tab:filter_hq} presents the best filter per cell line on HQ images. The key finding is that \textbf{no single filter is universally optimal}: DoG wins for 3/8 lines, Homomorphic for 2/8, and Butterworth/Elliptic/Gaussian each for 1/8. Improvement factors range from $+0.193$ (BT474) to $+0.474$ (SHSY5Y).

\begin{table}[H]
\centering
\caption{Best filter per cell line on HQ images ($n{=}808$). VF = improvement factor.}
\label{tab:filter_hq}
\begin{tabular}{@{}llrrrr@{}}
\toprule
\textbf{Cell Line} & \textbf{Best Filter} & \textbf{Best IoU} & \textbf{Raw IoU} & \textbf{$\Delta$IoU} & \textbf{VF} \\
\midrule
A172   & Homomorphic & 0.827 & 0.583 & +0.244 & 42\% \\
BT474  & Butterworth & 0.507 & 0.314 & +0.193 & 62\% \\
BV2    & DoG         & 0.527 & 0.191 & +0.336 & 176\% \\
Huh7   & Homomorphic & 0.609 & 0.285 & +0.324 & 114\% \\
MCF7   & Elliptic    & 0.639 & 0.338 & +0.302 & 89\% \\
SHSY5Y & DoG         & 0.799 & 0.325 & +0.474 & 146\% \\
SKOV3  & Gaussian    & 0.965 & 0.627 & +0.338 & 54\% \\
SkBr3  & DoG         & 0.630 & 0.183 & +0.447 & 244\% \\
\bottomrule
\end{tabular}
\end{table}

The filter$\times$cell-line heatmap (Fig.~\ref{fig:heatmap}) reveals structure: filters with smooth transitions (Gaussian, Butterworth) tend to perform consistently across lines, while sharp-cutoff filters (Ideal, Elliptic) show high variance.

% ---- Figure 2: Heatmap ----
\begin{figure}[H]
\centering
\includegraphics[width=\textwidth]{outputs/fig1_filter_matrix.pdf}
\caption{\textbf{Filter$\times$cell-line IoU heatmap.}
Mean IoU for 12 filter types across 8 cell lines on HQ images.
Color scale: blue = low IoU, red = high IoU. White cells indicate negative improvement.}
\label{fig:heatmap}
\end{figure}

% ---- 4.2 Filter Performance on LQ Images ----
\subsection{Filter improvements on low-quality images are 10--100$\times$ smaller}

Table~\ref{tab:filter_lq} presents filter performance on degraded images. The critical finding: filter improvements on LQ images are \textbf{10--100$\times$ smaller} than on HQ images.

\begin{table}[H]
\centering
\caption{Filter performance on degraded images. Transfer ratio = LQ improvement / HQ improvement.}
\label{tab:filter_lq}
\begin{tabular}{@{}llrrrrr@{}}
\toprule
\textbf{Degradation} & \textbf{Raw IoU} & \textbf{Best Filter} & \textbf{Best IoU} & \textbf{$\Delta$IoU} & \textbf{HQ $\Delta$IoU} & \textbf{Transfer} \\
\midrule
Noise $\sigma{=}50$ & 0.302 & Butterworth & 0.304 & +0.003 & +0.193 & \textbf{1.5\%} \\
Combined mild      & 0.281 & Butterworth & 0.309 & +0.028 & +0.210 & \textbf{13.3\%} \\
\bottomrule
\end{tabular}
\end{table}

% ---- Figure 3: Transfer Efficiency ----
\begin{figure}[H]
\centering
\includegraphics[width=0.85\textwidth]{outputs/fig4_transfer_efficiency.pdf}
\caption{\textbf{Filter transfer efficiency.} Transfer ratios from HQ to LQ for key filter types. All filters show poor transfer ($<$15\%).}
\label{fig:transfer}
\end{figure}

% ---- 4.3 Filter Transfer Efficiency ----
\subsection{Filter transfer from HQ to LQ is poor}

Transfer ratio (filter improvement on LQ / filter improvement on HQ) is consistently below 15\% for all filter types (Fig.~\ref{fig:transfer}). This means: (1) filter parameters optimized on HQ images do not transfer to LQ images; (2) different filter types may be needed for different quality levels; (3) the ``best'' filter depends on image quality, not just the cell line.

% ---- 4.4 Adaptive Filter Selection ----
\subsection{Adaptive selection improves mean IoU by $+0.130$}

Using cell-line-specific optimal parameters improved mean IoU from 0.378 (single fixed Butterworth) to 0.508 (best adaptive filter per line), a gain of $+0.130$ (Fig.~\ref{fig:adaptive}). The improvement is consistent across all cell lines.

% ---- Figure 4: Adaptive vs Fixed ----
\begin{figure}[H]
\centering
\includegraphics[width=\textwidth]{outputs/fig5_adaptive_vs_fixed.pdf}
\caption{\textbf{Adaptive vs fixed filter selection.}
\textbf{(a)} Mean IoU per cell line: fixed (single best filter) vs.\ adaptive (cell-line-specific optimal).
\textbf{(b)} IoU improvement distribution.}
\label{fig:adaptive}
\end{figure}

% ---- 4.5 Degradation-Specific Recommendations ----
\subsection{Degradation-specific filter recommendations}

Based on the complete evaluation, we provide degradation-specific recommendations (Table~\ref{tab:degradation_rec}). For noise-degraded images, Butterworth ($n{=}4$) with higher cutoff is optimal. For blur, DoG with tighter bandpass is preferred. For illumination shading, Homomorphic is the clear choice.

\begin{table}[H]
\centering
\caption{Degradation-specific filter recommendations.}
\label{tab:degradation_rec}
\begin{tabular}{@{}lllp{4cm}@{}}
\toprule
\textbf{Degradation} & \textbf{Best Filter} & \textbf{Key Parameters} & \textbf{Rationale} \\
\midrule
Gaussian noise    & Butterworth ($n{=}4$) & $d_{\text{low}}{=}0.03$, $d_{\text{high}}{=}0.25$ & Steeper roll-off removes noise \\
Defocus blur      & DoG                   & $\sigma_1{=}0.05$, $\sigma_2{=}0.20$ & Bandpass around cell frequency \\
Illumination      & Homomorphic           & $\gamma_L{=}0.3$, $\gamma_H{=}2.5$ & Designed for multiplicative artifacts \\
Combined mild     & Butterworth ($n{=}2$) & $d_{\text{low}}{=}0.02$, $d_{\text{high}}{=}0.30$ & Balanced approach \\
\bottomrule
\end{tabular}
\end{table}

% ---- 4.6 BBBC005 Validation ----
\subsection{Cross-dataset validation confirms quality-dependent selection}

On BBBC005 (25 blur levels), the best filter type varies with blur level (Fig.~\ref{fig:bbbc}): Butterworth dominates at low blur ($s_{01}$--$s_{10}$), DoG at high blur ($s_{16}$--$s_{25}$). The transition at $s_{12}$--$s_{14}$ confirms that quality-aware selection generalizes beyond LIVECell.

% ---- Figure 5: BBBC005 ----
\begin{figure}[H]
\centering
\includegraphics[width=0.85\textwidth]{outputs/ws2_blur_scale.png}
\caption{\textbf{BBBC005 blur-level analysis.} Filter recommendations across 25 blur levels. Transition from Butterworth to DoG at $s_{12}$--$s_{14}$.}
\label{fig:bbbc}
\end{figure}

% ---- 4.7 U-Net Segmentation Results ----
\subsection{U-Net segmentation confirms filter benefits}

With U-Net segmentation, bandpass preprocessing improves mean Dice from 0.72 (raw) to 0.79 (filtered), confirming that the benefits observed with Otsu thresholding extend to deep learning segmentation. Per-cell-line results are in Supplementary Table~S2.

% ---- Figure 6: U-Net Results ----
\begin{figure}[H]
\centering
\includegraphics[width=\textwidth]{outputs/report_fig6.png}
\caption{\textbf{Segmentation results.}
\textbf{(a)} U-Net Dice coefficient: raw vs.\ filtered, per cell line.
\textbf{(b)} Example segmentations: raw, filtered, ground truth.
\textbf{(c)} IoU improvement distribution across all images.}
\label{fig:unet_results}
\end{figure}

% ============================================================
%  5. DISCUSSION (~600 words)
% ============================================================
\section{Discussion}
\label{sec:discussion}

Our systematic evaluation yields four principal findings. First, \textbf{no single bandpass filter is universally optimal} for phase-contrast microscopy segmentation. The best filter depends on cell line morphology, with DoG excelling for 3/8 lines, Homomorphic for 2/8, and others for the remaining 3/8. This finding challenges the common practice of applying a single filter type to all images.

Second, \textbf{filter improvements on LQ images are 10--100$\times$ smaller} than on HQ images. This has critical implications for automated screening systems, where image quality is often degraded. Bandpass filtering alone cannot recover information lost during acquisition.

Third, \textbf{filter transfer from HQ to LQ is poor} ($<$15\%). Filter parameters must be re-optimized for each quality level, motivating quality-aware adaptive selection.

Fourth, \textbf{adaptive selection improves mean IoU by $+0.130$} over fixed filtering, a substantial improvement that is consistent across all cell lines.

\textbf{Comparison with related work.} Our results are consistent with the general image processing literature~\cite{gonzalez2018digital} recommending Butterworth for general-purpose filtering and Gaussian for artifact-free applications. Our contribution is the systematic evaluation across cell lines and degradation types, revealing the quality-dependence of filter performance. Previous microscopy studies compared at most 2--3 filter types on a single cell line; our 12-filter, 8-cell-line evaluation is the most comprehensive to date.

\textbf{Practical recommendations.} Based on our results, we recommend: (1) DoG ($\sigma_1{=}0.05$, $\sigma_2{=}0.20$) as the default for segmentation preprocessing; (2) Homomorphic ($\gamma_L{=}0.3$, $\gamma_H{=}2.5$) for images with illumination gradients; (3) Butterworth ($n{=}2$) as the safest general-purpose choice; (4) avoiding Ideal and Elliptic filters for biological images (ringing around cell boundaries); (5) implementing cell-line-adaptive filter selection for production systems.

\textbf{Limitations.} (1) Evaluation is limited to phase-contrast microscopy; results may not transfer to fluorescence or DIC. (2) Otsu thresholding, while standard, is a simple segmentation method; more sophisticated methods (Cellpose, StarConvex) may show different filter preferences. (3) The 25\% annotation subset limits some analyses. (4) Synthetic degradations may not capture all real-world quality issues.

% ============================================================
%  6. CONCLUSION (~200 words)
% ============================================================
\section{Conclusion}
\label{sec:conclusion}

We presented the first systematic evaluation of 12 bandpass filter types for phase-contrast microscopy segmentation across 8 cell lines and 13 degradation types, totaling over 20,000 segmentation evaluations. Key contributions:

\begin{enumerate}[nosep]
    \item No single filter is universally optimal---the best filter depends on cell line and image quality
    \item Filter improvements on LQ images are 10--100$\times$ smaller than on HQ images
    \item Filter transfer from HQ to LQ is consistently below 15\%
    \item Adaptive, cell-line-specific selection improves mean IoU by $+0.130$
    \item Evidence-based, degradation-specific filter recommendations
\end{enumerate}

All code, trained models, and the adaptive selection algorithm are available at \url{https://github.com/mjonyh/microscopic_images}.

% ============================================================
%  REFERENCES (~40)
% ============================================================
\bibliographystyle{plainnat}
\bibliography{references}

\end{document}
